\chapter{Elementos Pré-textuais}

Os elementos pré-textuais antecedem o conteúdo principal. Eles identificam o trabalho, registram informações institucionais e facilitam a navegação pelo documento. A capa não é contada. A contagem começa na folha de rosto, mas o número só aparece na primeira página textual. A página reservada aos dados internacionais de catalogação não participa da contagem.

\requisitonormativo

\section{Capa}

A capa identifica a instituição, o autor, o título, o local e o ano. O título deve representar com precisão o conteúdo do trabalho. Títulos excessivamente longos dificultam a leitura; títulos genéricos reduzem a identidade acadêmica do estudo.

No padrão institucional da UnDF, a capa deve valorizar a identidade da universidade e manter apresentação limpa. O curso não precisa aparecer na capa quando já constar na folha de rosto e na natureza do trabalho.

\section{Folha de rosto}

A folha de rosto apresenta dados essenciais para identificação acadêmica: autor, título, subtítulo quando houver, natureza do trabalho, curso, instituição, grau pretendido, orientação, local e ano.

A natureza do trabalho deve indicar que se trata de Trabalho de Conclusão de Curso apresentado ao curso correspondente, na Universidade do Distrito Federal Professor Jorge Amaury Maia Nunes -- UnDF, como requisito parcial para obtenção do grau definido.

\section{Dados internacionais de catalogação}

Os dados internacionais de catalogação devem seguir o código de catalogação vigente e o fluxo definido pela Biblioteca Central da UnDF. A ABNT NBR 14724:2024 determina sua posição no verso da folha de rosto ou, em documento eletrônico, na página imediatamente posterior, mas a exigência, a emissão e o momento de solicitação dependem da orientação institucional.

O estudante não deve criar sozinho uma ficha catalográfica nem inventar classificação, notação de autor ou descritores. A página que contém esses dados não é contada nem numerada.

\fluxopendente

Devem ser confirmados pela Biblioteca Central da UnDF o serviço responsável, os critérios de emissão, o momento da solicitação, o formato aceito e a relação com o depósito no repositório.

\section{Errata}

A errata é elemento eventual. Ela registra correções identificadas depois da finalização ou impressão do trabalho. Deve conter a referência do trabalho e a indicação objetiva de folha, linha, trecho incorreto e forma correta.

Não se deve usar errata para corrigir problemas extensos de redação, estrutura ou normalização que poderiam ter sido resolvidos antes da entrega.

\section{Folha de aprovação}

A folha de aprovação registra a avaliação do trabalho, a data e os membros da banca. A ABNT exige nome, titulação acadêmica e instituição dos integrantes. No padrão institucional validado no âmbito do projeto, a titulação aparece de forma sintética antes do nome, por exemplo, ``Prof. Dr. Nome'', sem uma segunda linha de título acadêmico. A data de aprovação segue a forma: Brasília, DF, dia de mês de ano.

Os nomes da banca devem ser apresentados de forma uniforme, com prefixo acadêmico antes do nome. A titulação não deve ser repetida de maneira extensa em linha separada.

\decisaoinstitucional

\section{Dedicatória, agradecimentos e epígrafe}

Dedicatória, agradecimentos e epígrafe são opcionais. A dedicatória é breve e pessoal. Os agradecimentos podem reconhecer orientação, apoio institucional, colaboração técnica, financiamento, familiares e pessoas que contribuíram para o percurso.

A epígrafe pré-textual é uma citação relacionada ao tema geral do trabalho, apresentada conforme o modelo gráfico institucional. Deve indicar autoria e possuir referência correspondente. A epígrafe também pode aparecer na abertura de uma seção primária; nesse caso, relaciona-se especificamente ao conteúdo da seção e continua sujeita às regras de citação e referência.

Esses elementos devem manter sobriedade acadêmica. Textos muito longos, excessivamente informais ou desconectados do contexto institucional devem ser evitados.

\section{Resumo}

O resumo apresenta uma síntese do trabalho em um único parágrafo. Deve indicar objetivo, método, resultados principais e conclusão ou considerações finais, em linguagem clara e concisa. Para trabalhos acadêmicos, a ABNT NBR 6028:2021 recomenda de 150 a 500 palavras. Ele não deve conter citações, notas, figuras, tabelas, enumerações extensas ou informações que não estejam desenvolvidas no corpo do trabalho.

As palavras-chave devem representar conceitos centrais do trabalho. Devem aparecer logo abaixo do resumo, após a expressão ``Palavras-chave:'', separadas por ponto e vírgula e encerradas por ponto. A ABNT não estabelece uma quantidade universal; use quantidade adequada ao conteúdo e escolha termos específicos que favoreçam a recuperação da informação.

\decisaoinstitucional

\section{Resumo em língua estrangeira}

O resumo em língua estrangeira corresponde ao resumo em português, vertido para outro idioma. O idioma mais comum é o inglês, mas o curso ou a instituição pode admitir outro idioma conforme a área.

O resumo estrangeiro deve preservar equivalência de conteúdo, não ser uma versão reduzida, ampliada ou divergente do resumo em português.

\caixaatencao{Resumo em português e resumo em língua estrangeira são elementos obrigatórios distintos na ABNT NBR 14724:2024. A equivalência de conteúdo não autoriza a supressão de um deles.}

\section{Listas}

Listas de ilustrações, tabelas, abreviaturas, siglas e símbolos são usadas quando ajudam o leitor a localizar itens recorrentes. Não é necessário criar lista para quantidade mínima de elementos quando isso não favorece a consulta.

\section{Sumário}

O sumário apresenta a estrutura interna do trabalho e deve refletir fielmente os títulos e páginas do texto. Elementos pré-textuais não entram no sumário. O sumário é o último elemento pré-textual.

Em arquivos digitais, recomenda-se que as entradas sejam links navegáveis. O sumário deve ser atualizado somente depois da paginação final e conferido no PDF entregue.

\section{Exemplo de contagem e exibição}

O Apêndice \ref{ap:paginacao} apresenta uma sequência completa. Em síntese: a capa fica fora da contagem; a folha de rosto inicia a contagem; a página dos dados catalográficos é desconsiderada; os demais pré-textuais continuam a contagem sem mostrar o número; e a primeira página da introdução exibe o número correspondente à sua posição.
